\documentclass[a4paper, serif, 10pt]{moderncv}

%% moderncv
% color and style
\moderncvcolor{black}
\moderncvstyle{banking}

% renew moderncv command to adapt for CJK colon
\renewcommand*{\cvitem}[3][.25em]{%
  \ifstrempty{#2}{}{\hintstyle{#2}：}{#3}%
  \par\addvspace{#1}}

\renewcommand*{\cvitemwithcomment}[4][.25em]{%
  \savebox{\cvitemwithcommentmainbox}{\ifstrempty{#2}{}{\hintstyle{#2}：}#3}%
  \setlength{\cvitemwithcommentmainlength}{\widthof{\usebox{\cvitemwithcommentmainbox}}}%
  \setlength{\cvitemwithcommentcommentlength}{\maincolumnwidth-\separatorcolumnwidth-\cvitemwithcommentmainlength}%
  \begin{minipage}[t]{\cvitemwithcommentmainlength}\usebox{\cvitemwithcommentmainbox}\end{minipage}%
  \hfill% fill of \separatorcolumnwidth
  \begin{minipage}[t]{\cvitemwithcommentcommentlength}\raggedleft\small\itshape#4\end{minipage}%
  \par\addvspace{#1}}

%% page layout/margins
\usepackage[top=2.5cm, bottom=2.5cm, left=1.5cm, right=1.5cm]{geometry}

\nopagenumbers{}

%% Babel config for English language
\usepackage[english]{babel}

%% fontspec
\usepackage{fontspec}

\IfFontExistsTF{Linux Libertine}{
  \setmainfont[Ligatures={TeX, Common}, Numbers=Lining, ItalicFont=Linux Libertine]{Linux Libertine}
}{}
\IfFontExistsTF{Linux Libertine O}{
  \setmainfont[Ligatures={TeX, Common}, Numbers=Lining, ItalicFont=Linux Libertine O]{Linux Libertine O}
}{}

%% CTeX
% CJK support, used to show CJK characters in the resume
%
% - heading=false: disable ctex heading
% - punct=kaiming: use kaiming punctuations styles for CJK
% - scheme=plain: use plain scheme, do not override \normalsize font size
% - space=auto: space settings for CJK characters
%
% CJK fonts are resolved per locale language, e.g. Noto Serif/Sans CJK JP
% for Japanese resumes, with fallback to the next region if not installed
%
% ref:
% - http://ctan.mirrorcatalogs.com/language/chinese/ctex/ctex.pdf
\usepackage[UTF8, heading=false, punct=kaiming, scheme=plain, space=auto]{ctex}

\IfFontExistsTF{Noto Serif CJK JP}{
  \setCJKmainfont{Noto Serif CJK JP}
}{}
\IfFontExistsTF{Noto Sans CJK JP}{
  \setCJKsansfont{Noto Sans CJK JP}
}{}

%% line spacing
\usepackage{setspace}
\setstretch{1.125}

%% URL styling - use normal text instead of monospace
\urlstyle{same}

%% Patch \httplink and \httpslink to handle full URLs
% The original moderncv macros always prepend http:// or https:// to URLs.
% This causes issues when the URL already has a protocol (e.g., https://example.com)
% resulting in malformed URLs like https://https://example.com.
% This patch checks if the URL already contains "://" and uses it directly if so.
% Note: We use \str_set:Nx (x-expansion) to expand macros like \@homepage first.
\ExplSyntaxOn
\str_new:N \g_yamlresume_url_str

\RenewDocumentCommand{\httpslink}{O{}m}{
  \str_set:Nx \g_yamlresume_url_str {#2}
  \str_if_in:NnTF \g_yamlresume_url_str {://}
    { \url{#2} }
    { \url{https://#2} }
}

\RenewDocumentCommand{\httplink}{O{}m}{
  \str_set:Nx \g_yamlresume_url_str {#2}
  \str_if_in:NnTF \g_yamlresume_url_str {://}
    { \url{#2} }
    { \url{http://#2} }
}
\ExplSyntaxOff

\name{佐藤 美咲}{}
\title{人事マネージャー | 組織開発・人材育成の専門家}
\phone[mobile]{+81 90-1234-5678}
\email{misaki.sato@example.com}
\homepage{https://www.misaki-sato-hr.com}

\address{日本、東京、東京都、千代田区千代田1-1、100-0001}{}{}

\extrainfo{{\small \faLinkedin}\ \href{https://www.linkedin.com/in/misaki-sato-hr}{@misaki-sato-hr} {} {} {} • {} {} {} 
{\small \faTwitter}\ \href{https://twitter.com/misaki\_hr}{@misaki\_hr}}

\begin{document}

\maketitle

\section{基本情報}

\cvline{}{\begin{itemize}
\item 人事歴10年以上の経験を持つ人事マネージャー。採用から育成、評価制度の設計・運用まで幅広く担当。
\item 組織開発とエンゲージメント向上に注力し、データドリブンな人事施策を推進。
\item グローバル環境での人事業務経験あり。英語でのコミュニケーションも可能。
\item 労働法や社会保険に関する深い知識を持ち、労務リスク管理にも強み。
\end{itemize}}
\section{学歴}

\cventry{2010年4月～2014年3月}
        {学士、経営学部 マネジメント学科、成績：3.8/4.0}
        {青山学院大学}
        {\url{https://www.aoyama.ac.jp}}
        {}
        {\begin{itemize}
\item 人事管理や組織行動学を中心に履修し、理論と実践の両面から人材マネジメントを学ぶ
\item ゼミでは企業の人事制度に関する研究を行い、卒業論文を執筆
\item 学内のキャリア支援ボランティアに参加し、学生の就職活動をサポート
\end{itemize}
\textbf{コース}：人事管理論、組織行動学、労働法、産業心理学、統計学、マーケティング}
\section{職歴}

\cventry{2020年1月～現在}
        {人事マネージャー}
        {株式会社テクノロジーソリューションズ}
        {\url{https://www.tech-sol.co.jp}}
        {}
        {\begin{itemize}
\item 人事部門のマネージャーとして、採用・育成・評価・労務の全般を統括
\item 年間50名以上の採用をリードし、採用プロセスの効率化と候補者体験の向上を実現
\item 新評価制度（OKR＋360度評価）を設計・導入し、組織のパフォーマンスを向上
\item 従業員エンゲージメント調査を実施し、離職率を15\%低減
\item 労務管理のデジタル化を推進し、業務効率を30\%改善
\end{itemize}
\textbf{キーワード}：採用、育成、評価制度、労務管理、組織開発}

\cventry{2014年4月～2019年12月}
        {人事スペシャリスト}
        {株式会社グローバル商事}
        {\url{https://www.global-shoji.co.jp}}
        {}
        {\begin{itemize}
\item 新卒・中途採用を担当し、年間30名以上の採用を実現
\item 研修プログラムの企画・運営を行い、新入社員の早期戦力化を支援
\item 給与計算・社会保険手続きなどの労務業務を正確かつ迅速に処理
\item 人事データの分析を行い、経営層へのレポーティングを実施
\end{itemize}
\textbf{キーワード}：採用、研修、給与計算、社会保険}
\section{言語}

\cvline{日本語}{ネイティブ・母国語レベル \hfill \textbf{キーワード}：漢字検定1級、日本語教育能力検定}
\cvline{英語}{完全な実務レベル \hfill \textbf{キーワード}：TOEIC 950、英検1級}
\section{スキル}

\cvline{人事戦略}{エキスパート \hfill \textbf{キーワード}：採用計画、人材育成、組織設計、パフォーマンス管理}
\cvline{労務管理}{上級 \hfill \textbf{キーワード}：労働法、社会保険、給与計算、就業規則}
\cvline{データ分析}{中級 \hfill \textbf{キーワード}：Excel、HR分析、統計}
\cvline{コミュニケーション}{エキスパート \hfill \textbf{キーワード}：交渉、プレゼンテーション、ファシリテーション}
\section{受賞・表彰}

\cventry{2022年12月}
        {株式会社テクノロジーソリューションズ}
        {年間最優秀人事賞}
        {}
        {}
        {人事改革プロジェクトのリーダーとして、採用・育成・評価制度の刷新に貢献し、組織のエンゲージメントスコアを20\%向上させた。}
\section{資格・免状}

\cventry{2018年8月}
        {厚生労働省}
        {社会保険労務士}
        {\url{https://www.mhlw.go.jp}}
        {}
        {}

\cventry{2019年11月}
        {厚生労働省}
        {キャリアコンサルタント}
        {\url{https://www.mhlw.go.jp}}
        {}
        {}
\section{出版・著作}

\cventry{2023年6月}
        {次世代人事戦略：データドリブンな人材マネジメント}
        {日本の人事部}
        {\url{https://www.jinji-bu.jp}}
        {}
        {人事データ分析を活用した採用・育成の最適化について寄稿。データに基づく意思決定の重要性と具体的な手法を解説。}
\section{プロジェクト}

\cventry{2022年4月～2023年3月}
        {プロジェクトの詳細説明。評価制度の設計からシステム選定、運用開始までを一貫して担当。}
        {人事評価システム刷新プロジェクト}
        {\url{https://www.tech-sol.co.jp/projects/hr-system}}
        {}
        {評価制度の見直しとシステム導入をリード。360度評価とOKRを組み合わせた新評価制度を設計し、クラウド人事システムを導入。
\textbf{キーワード}：評価制度、人事システム、プロジェクトマネジメント}
\section{興味・関心}

\cvline{読書}{ビジネス書、心理学、歴史小説}
\cvline{旅行}{国内旅行、海外旅行、温泉}
\cvline{ヨガ}{ヨガインストラクター資格、瞑想}
\section{ボランティア活動}

\cventry{2019年4月～現在}
        {キャリアアドバイザー}
        {NPO法人 若者キャリア支援協会}
        {\url{https://www.example-npo.org}}
        {}
        {大学生や若年求職者向けのキャリア相談を実施。履歴書添削や面接対策を通じて、就職活動を支援。}
    
\end{document}